% Part: second-order-logic
% Chapter: metatheory
% Section: second-order-arithmetic

\documentclass[../../../include/open-logic-section]{subfiles}

\begin{document}

\olfileid[fr]{sol}{met}{spa}

\olsection{L'arithmétique du second ordre}

Rappelons que la théorie $\Th{PA}$ de l'arithmétique
de Peano comprend les huit axiomes de~$\Th{Q}$ :
\begin{align*}
& \lforall[x][\eq/[x'][\Obj 0]]\\
& \lforall[x][\lforall[y][(\eq[x'][y'] \lif \eq[x][y])]]\\
& \lforall[x][(\eq[x][\Obj 0] \lor \lexists[y][\eq[x][y']])]\\
& \lforall[x][\eq[(x + \Obj 0)][x]]\\
& \lforall[x][\lforall[y][\eq[(x + y')][(x + y)']]]\\
& \lforall[x][\eq[(x \times \Obj 0)][\Obj 0]]\\
& \lforall[x][\lforall[y][\eq[(x \times y')][((x \times y) + x)]]]\\
& \lforall[x][\lforall[y][(x < y \liff \lexists[z][\eq[(z' + x)][y]])]]\\
\intertext{ainsi que la clôture universelle sur les paramètres libres de chaque instance de la forme}
& (!A(\Obj 0) \land \lforall[x][(!A(x) \lif !A(x'))]) \lif \lforall[x][!A(x)].
\end{align*}
Cette dernière expression est un \emph{schéma} :
elle engendre une infinité d'!!{sentence}s du
langage arithmétique, un pour chaque
!!{formula}~$!A(x)$. On l'appelle le
\emph{schéma d'axiomes d'induction} du premier
ordre.
Dans l'arithmétique de Peano \emph{du second ordre}
$\Th{PA^2}$, l'induction peut s'exprimer par
un seul !!{sentence}. Cette théorie comprend
les huit axiomes ci-dessus et l'\emph{axiome
d'induction} du second ordre :
\[
\lforall[X][((X(\Obj 0) \land \lforall[x][(X(x) \lif X(x'))]) \lif \lforall[x][X(x)])].
\]
L'axiome dit que tout sous-ensemble~$X$ du
!!{domain} qui contient $\Assign{\Obj{0}}{M}$
et qui, avec chaque $x\in\Domain{M}$,
contient aussi $\Assign{\prime}{M}(x)$
(autrement dit, qui est \emph{clos par successeur})
contient tout le !!{domain} : $X=\Domain{M}$.

Sous la sémantique standard, cet axiome garantit
qu'une !!{structure} qui le satisfait ne possède
que les !!{element}s imposés par les axiomes,
c'est-à-dire les valeurs des numéraux~$\num n$
pour $n\in\Nat$. Un modèle de~$\Th{PA^2}$
ne contient aucun nombre non standard.

\begin{thm}
\ollabel{thm:sol-pa-standard}
Si $\Sat{M}{\Th{PA^2}}$, alors
$\Domain{M}=\Setabs{\Value{\num{n}}{M}}{n\in\Nat}$.
\end{thm}

\begin{proof}
Posons
$N=\Setabs{\Value{\num{n}}{M}}{n\in\Nat}$
et supposons $\Sat{M}{\Th{PA^2}}$. Pour chaque
$n\in\Nat$, la valeur $\Value{\num{n}}{M}$
appartient au !!{domain}; donc
$N\subseteq\Domain{M}$.

Montrons l'inclusion inverse. Considérons une
assignation~$s$ telle que $s(X)=N$. Par hypothèse,
\begin{align*}
\Sat{M}{\lforall[X][((X(\Obj 0) \land \lforall[x][(X(x)
      \lif X(x'))]) \lif \lforall[x][X(x)])]}, & \text{ donc}\\
\Sat{M}{(X(\Obj 0) \land \lforall[x][(X(x) \lif X(x'))]) \lif
  \lforall[x][X(x)]}[s]. &
\end{align*}
L'antécédent de cette implication est satisfait.
En effet, $\Value{\Obj{0}}{M}\in N$, d'où
$\Sat{M}{X(\Obj{0})}[s]$. Le second conjoint,
$\lforall[x][(X(x)\lif X(x'))]$, l'est aussi :
si $x\in N$, alors, pour un certain~$n$,
$x=\Value{\num{n}}{M}$, et par conséquent
$\Assign{\prime}{M}(x)=
\Value{\num{n+1}}{M}\in N$.
Nous avons donc
$\Sat{M}{X(\Obj 0) \land
\lforall[x][(X(x) \lif X(x'))]}[s]$,
puis $\Sat{M}{\lforall[x][X(x)]}[s]$.
Ainsi, tout $x\in\Domain{M}$ appartient à
$s(X)=N$, ce qui donne $\Domain{M}\subseteq N$.
\end{proof}

\begin{cor}
\ollabel{cor:sol-pa-categorical}
Deux modèles quelconques de~$\Th{PA^2}$ sont
isomorphes.
\end{cor}

\begin{proof}
D'après le \olref{thm:sol-pa-standard}, le !!{domain}
de tout modèle de~$\Th{PA^2}$ est épuisé par
les valeurs $\Value{\num{n}}{M}$. Un tel modèle
satisfait aussi~$\Th{Q}$. Par la
\olref[mod][mar][stm]{prop:thq-standard}, il est
standard, donc isomorphe à~$\Struct{N}$.
\end{proof}

Nous avons défini $\Th{PA^2}$ au moyen des huit
axiomes arithmétiques et de l'axiome d'induction
du second ordre. Grâce à la puissance expressive
de cette logique, les \emph{deux premiers} axiomes
arithmétiques, joints à l'induction, suffisent
en réalité pour l'arithmétique de Peano du second
ordre.

\begin{prop}\ollabel{prop:sol-pa-definable}
Soit $\Th{PA^{2\dagger}}$ la théorie du second ordre
dans le langage réduit au zéro et au successeur, constituée des deux premiers axiomes arithmétiques
(ceux du successeur) et de l'axiome d'induction
du second ordre. Les relations et opérations
$\le$, $+$ et $\times$ y sont définissables.
\end{prop}

\begin{proof}
Pour définir $\le$, cherchons une !!{formula}
$!A_\le(x,y)$ telle que
$\Sat{N}{!A_\le(\num n,\num m)}$ si et seulement
si $n\le m$. Considérons
\[
!B(x,Y) \ident Y(x) \land \lforall[y][(Y(y) \lif Y(y'))].
\]
Tout ensemble $Y\subseteq\Nat$ qui satisfait~$!B(\num n,Y)$
contient $\Setabs{m}{n\le m}$; réciproquement,
ce dernier ensemble satisfait lui-même cette !!{formula}.
Nous pouvons donc poser
$!A_\le(x,y)\ident
\lforall[Y][(!B(x,Y)\lif Y(y))]$.

Pour l'addition, observons que $k+l=m$ si et
seulement s'il existe une fonction~$u$ telle que
$u(0)=k$, $u(n')=u(n)'$ pour tout~$n$ et
$u(l)=m$. Cette équivalence permet de définir
l'addition dans~$\Th{PA^{2\dagger}}$ par :
\[
!A_+(x,y,z) \ident \lexists[u][(u(\Obj 0)=x \land
\lforall[w][u(w')=u(w)'] \land u(y)=z)].
\]
Il en résulte que
$\Sat{N}{!A_+(\num k,\num l,\num m)}$
si et seulement si $k+l=m$.
\end{proof}

\begin{editorial}
La formule de l'addition dans la source quantifie
$w$ mais écrit $u(x')=u(x)'$ dans la clause
récursive. Nous y remplaçons $x$ par $w$ : la
clause exprime alors bien la récurrence pour
tout~$w$. Les conclusions de catégoricité et
de définissabilité utilisent la sémantique
standard du second ordre, où $X$ parcourt tous
les sous-ensembles et $u$ toutes les fonctions
du domaine concerné.
\end{editorial}

\begin{prob}
Complétez la preuve de la
\olref[sol][met][spa]{prop:sol-pa-definable}.
\end{prob}

\end{document}
